\cvsection{Expériences professionnelles}

\begin{cventries}
  \cventry
    {Stagiaire ingénieur de recherche en IA}
    {INRIA}
    {Bordeaux, France}
    {Mars 2026 -- Septembre 2026}
    {
      \begin{cvitems}
        \item {Conception en autonomie, avec le suivi de mes encadrants, d’une application et d’un agent IA multi-étapes répondant au besoin d’aider les chercheurs à trouver des articles scientifiques.}
        \item {Développement de l’agent en Python sans framework d’orchestration pour comprendre ses mécanismes : intégration de l’API OpenAI, appels d’outils, enchaînement des étapes, gestion de l’état et de la mémoire.}
        \item {Conception de prompts adaptés au cas d’usage et à l’utilisation des outils, avec gestion du contexte et sélection par l’utilisateur des messages à transmettre au LLM.}
        \item {Automatisation de la collecte de métadonnées d’articles scientifiques en ligne et transmission des résultats de tâches cron par webhooks.}
        \item {Développement d’un backend Python/FastAPI avec API REST documentées avec OpenAPI, persistance MongoDB et interface React/Next.js pour interagir avec l’agent et ses outils.}
        \item {Mise en œuvre d’un pipeline RAG avec chunking, embeddings et recherche sémantique dans une base vectorielle ChromaDB.}
        \item {Conteneurisation des composants avec Docker et mise en place d’une CI/CD sur GitHub.}
        \item {Évaluation d’OpenClaw dans un environnement Docker isolé et analyse des risques liés à l’accès autonome de l’agent aux outils et aux fichiers.}
        \item {Gestion d’une application utilisant des LLM installés localement sur un Mac Studio et utilisation de vLLM pour servir des modèles sur une infrastructure HPC équipée de GPU NVIDIA.}
        \item {Comparaison de Qwen, Gemma, GLM et de plusieurs quantifications sur HPC avec Slurm selon la qualité des réponses, la latence, le TTFT et la taille du contexte.}
        \item {Définition du protocole expérimental, analyse des résultats et rédaction en cours d’un article scientifique comparant les modèles.}
      \end{cvitems}
    }

  \cventry
    {Stagiaire de recherche en apprentissage par renforcement}
    {Université Karunya}
    {Coimbatore, Inde}
    {Juin 2025 -- Sept. 2025}
    {
      \begin{cvitems}
        \item {Développement d’un modèle d’apprentissage par renforcement et de son environnement d’entraînement pour la navigation autonome d’un robot hospitalier.}
        \item {Déploiement et optimisation du logiciel et du modèle sur Raspberry Pi pour des tests en conditions réelles.}
      \end{cvitems}
    }
\end{cventries}
